\documentclass[12pt]{report}

\usepackage[utf8]{inputenc}
\usepackage[french]{babel}
\usepackage{fullpage}
\usepackage{graphicx}
\usepackage{fancyhdr}	% headers/footers
\usepackage{xcolor}		% to use our own color
\usepackage{lastpage}	% to easily know the total number of pages
\usepackage{titling}	% to easily know the total number of pages
\usepackage{colortbl}	% to put color in a table background
\usepackage{datetime}	% to allow us set a new date formatting

% Custom defines zone

% Define useful hand-made colors
\definecolor{epiBlue}{RGB}{23,54,93}
\definecolor{lightGray}{gray}{0.85}

% Bit of code to bold an entire table row
% http://tex.stackexchange.com/questions/4811/make-first-row-of-table-all-bold
\newcolumntype{$}{>{\global\let\currentrowstyle\relax}}
\newcolumntype{^}{>{\currentrowstyle}}
\newcommand{\rowstyle}[1]{\gdef\currentrowstyle{#1}%
  #1\ignorespaces
}

% Defining a "dd/mm/yyyy" date format
\newdateformat{dashDate}{\twodigit{\THEDAY}/\twodigit{\THEMONTH}/\twodigit{\THEYEAR}}

% Define Document Title
\newcommand{\DocTitle}{[2015][UbuntuOne][EDE]}

% end of Defines


% fancyhdr-specific commands
\setlength{\headheight}{15.2pt}

%% Defining headers and footers contents.

% Big dirty hack of the "empty" pagestyle to show header and footer on the title page (in wait of a better solution)
\fancypagestyle{empty}
{
	\renewcommand{\headrulewidth}{0pt}
	\renewcommand{\footrulewidth}{1pt}
	\fancyhead[L]{\includegraphics[height=42pt]{../logo_eip.png}
	}
	\fancyhead[R]{\colorbox{epiBlue}{\color{white}\textbf{\Large{\DocTitle} } }
	}

	\fancyfoot[L]{
   		\textcolor{gray}{\{EPITECH.\}}
	}
	\fancyfoot[C]{
		\jobname
	}
	\fancyfoot[R]{}
}

\fancypagestyle{EIP}
{
	\renewcommand{\headrulewidth}{0pt}
	\renewcommand{\footrulewidth}{1pt}
	\fancyhead[L]{\includegraphics[height=42pt]{../logo_eip.png}
	}
	\fancyhead[R]{\colorbox{epiBlue}{\color{white}\textbf{\Large{\DocTitle} } }
	}

	\fancyfoot[L]{
   		\textcolor{gray}{\{EPITECH.\}}
	}
	\fancyfoot[C]{
		\jobname
	}
	\fancyfoot[R]{
		\thepage/\pageref{LastPage}
	}
}

\pagestyle{EIP} % does not seem to work ...

% end of fancyhdr stuff

%Gummi|063|=)

%\title{The Title\\\normalsize A Sub-title}
\title{
	\huge{\textbf{\textcolor{epiBlue}{EIP Ubuntu One} } }\\
	\Large{\textbf{\emph{\textcolor{gray}{Étude de l'existant (EDE)} } } }
}


\begin{document}
\addtocontents{toc}{\protect\refstepcounter{page}} % makes the table of contents count pages from 1 (one)
\maketitle

\thispagestyle{empty}
\vspace*{10mm}

\textbf{\emph{\textcolor{onitu}{\large{Résumé du document} } } }\\

Résumé !!!

\clearpage


\thispagestyle{empty}
\vspace*{10mm}
\textbf{\emph{\textcolor{epiBlue}{Description du document} } }\\

\begin{tabular}{|>{\columncolor[gray]{0.85}\color{epiBlue} \bfseries } l|l|}
\hline
	Titre & \DocTitle\\
\hline
	Date & \dashDate\today \\
\hline
	Auteur & Alexandre BARON\\
\hline
	Responsable & Louis Roché\\
\hline
	E-Mail & serveurubintu\_un2015@labeip.epitech.eu\\
\hline
	Sujet & Étude de l'existant autour de notre projet\\
\hline
	Mots clés & Étude, état de l'art, cloud\\
\hline
	Version du modèle & 1.0\\
\hline
\end{tabular}
\vspace*{10mm}

\textbf{\emph{\textcolor{epiBlue}{Tableau des révisions} } }\\


\begin{tabular}{|$l|p{4cm}|p{2cm}|p{5cm}|}
\hline
\rowcolor{lightGray}
\rowstyle{ \color{epiBlue} \bfseries}
	Date & Auteur & Section(s) & Commentaires\\
\hline
	05/02/2013 & Alexandre BARON & Toutes & Première version \\
\hline
	12/02/2013 & Alexandre BARON & Chapitre 1 & Remplissage du Rappel sur l'EIP \\
\hline
	12/02/2013 & Wannes ROMBOUTS & Chapitre 2 & Ajout de l'étude sur Dropbox \\
\hline
	12/02/2013 & Maxime CONSTANTINIAN & Chapitre 3 & Ajout de l'étude sur SparkleShare \\
\hline
	12/02/1013 & Antoine ROZO & Chapitre 4 & Ajout de l'étude sur Google Drive, iCloud et SkyDrive \\
\hline
	12/02/2013 & Louis ROCHÉ & Chapitre 5 & Ajout de l'étude sur Syncany \\
\hline
	12/02/2013 & Yannick PÉROUX & Chapitre 6 & Ajout de l'étude sur OwnCloud \\
\hline
	12/02/2013 & Louis ROCHÉ & Chapitres 5 et 6 & Corrections mineures \\
\hline
	12/02/2013 & Yannick PÉROUX & Chapitre 3 & Améliorations de l'étude sur SparkleShare \\
\hline
	& & & \\
\hline
\end{tabular}

\tableofcontents
\thispagestyle{empty}

\chapter{Rappel de l'EIP}
\setcounter{page}{1} %reset the page count

\section{Qu'est-ce qu'un EIP et Epitech}
Epitech, école d'expertise informatique en cinq ans, propose aux étudiants, à partir de leur troisième année, un projet de fin d'études: l'EIP (pour \emph{Epitech Innovative Project}).\\

À ce titre, les élèves doivent s'organiser en un groupe d'au moins cinq personnes et choisir un sujet porteur de nouveautés ou améliorant un ancien sujet. L'EIP est un passage obligatoire et unique dans la scolarité de l'étudiant, de par son envergure (18 mois) et la préparation requise. Le but est, à la fin du temps imparti, d'obtenir un projet commercialisable.


\section{Sujet de notre EIP}
    Il s’agit d’un projet visant à proposer une implémentation libre et Open Source du serveur d’Ubuntu One.\\

    Ubuntu One est un service de Canonical (sponsor officiel d'Ubuntu) permettant de disposer d’un espace de stockage en ligne qui sera synchronisé entre différents ordinateurs et périphériques compatibles via un logiciel client. Le client et le protocole d’Ubuntu One sont disponibles sous licence libre. Néanmoins, le serveur est propriétaire et n’a pas été publié.\\


    Notre objectif est donc de proposer un équivalent libre de ce serveur, afin
    de profiter des fonctionnalités d’Ubuntu One par exemple à des fins d’autohébergement.\\

    Le client officiel n’étant pas capable d’utiliser un serveur différent de celui
    de Canonical, nous prévoyons également d'effectuer un fork afin d'ajouter cette option.

\thispagestyle{EIP} % seems mandatory

\chapter{Solutions existantes}

\chapter{Dropbox}
\thispagestyle{EIP} % seems mandatory
\section{Présentation}
DropBox est un service de stockage dans le cloud qui permet la syncronisation de fichiers entre différants terminaux. Dropbox à longtemps été la solution de référence et est énormément utilisé. Il est compatible avec GNU/Linux, Windows, Mac, Blackberrry, iOS et Android.\\

\section{Historique}
Le projet Dropbox est né au MIT en 2007 et à été lancé un an plus tard. En 2011 OPSWAT rapporte que Dropbox represente 14.14\% du marché mondial et ette même année il dépasse les 50 million d'utilisateurs. En 2012 ce chiffre est doublé et Dropbox annonce 100 million d'utilisateurs.\\

\section{Description}
Dropbox crée un dossier spécial sur chaque ordinateur ou il est installé. Il va syncroniser ce dossier entre les diférents terminaux en répercutant les modifications apportées aux fichiers ou sous-dossiers. Les données placés dans ce dossier sont aussi accesibles depuis une interface web.\\

Les utilisateurs de dropbox ont gratuitement acces à un espace de stockage de 18Go, avec la possibilitée de prendre un compte Pro pour bénéficier d'espace suplémentaire en payant un abonement mensuel.\\

D'un point de vue technique le server et le client Dropbox sont tout deux écrit en Python en utilisant des librairies standarts tel que Twisted et ctypes. La gestion de l'historique d'un fichier est similaire à celle d'un gestionaire de version clasique dans le sens ou il enregistre uniquement les diférences entre deux version succesive d'un fichier, c'est le delta encoding.\\

Dropbox utilise le Amazon S3 pour ses serveurs.\\

\section{Critiques}

En juillet 2011, Dropbox a modifié ses conditions d'utilisation et peut maintenant utiliser les fichiers stockés par ses clients sans leur autorisation. Ceci à poussé de nombreux utilisateurs à abondonner le service et a s'orienter vers d'autres solutions ou l'auto-hébergement.\\

\chapter{Sparkleshare}
\thispagestyle{EIP} % seems mandatory
\section{Présentation}
SparkleShare est une solution de stockage dans le cloud, orienté collaboration. Il repose sur Git pour gérer la synchronisation des fichiers. C'est un logiciel libre (GPLv3), compatible GNU/Linux, Windows, Mac, iOS et Android.\\

\section{Historique}
La version 1.0 est sortie le 9 décembre 2012 mais le travail a commencé début 2011 avec la première version publique le 14 fevrier 2011.\\

\section{Description}
Codé en C\#, SparkleShare repose essentiellement sur Git. Il est compatible avec n'importe quel dépôt Git, et peut donc être utilisé avec des services comme Github, Bitbucket ou Gitorious.\\

Son utilisation est essentiellement destinée aux utilisateurs qui voudraient travailler à plusieurs sur les mêmes fichiers, de manière plus conviviale qu'un outil de gestion de version classique.\\

\section{Critiques}
Aujourd'hui, peu d'utilisateurs semblent avoir adopté Sparkleshare, ce qui est probablement dû à son installation relativement complexe, ainsi qu'à son utilité limitée par rapport à des solutions comme Dropbox ou OwnCloud.\\

Git fonctionne très bien avec des fichiers textes, mais a plus de mal avec des binaires. Il est donc déconseillé d'utiliser Sparkleshare avec des fichiers multimédias.\\
Qui plus est, Sparkleshare ne propose pas de compatibilité avec les protocoles plus classiques, comme HTTP, FTP ou Webdav.\\
\chapter{Services de cloud par navigateur}
\thispagestyle{EIP} % seems mandatory

\section{Google drive}

\subsection{Présentation}
Google Documents est un service web créé en 2006, regroupant un tableur, un traitement de textes et un éditeur de diaporamas.
Lancé en 2012, Google Drive lui succède, proposant en plus un système de stockage en ligne, valable pour un grand nombre de types de fichiers.\\

\subsection{Description}
DONNÉES BRUTES\\
Fonctionnalités\\
* Ce système permet le partage de données avec d'autres membres, mais aussi, par son interface, une édition simultanée d'un même document par plusieurs personnes (travail collaboratif)\\
* Streaming de fichiers multimédia\\
* Permet de ne synchroniser que certains dossiers préalablement sélectionnés\\
* Intégration Gmail, G+\\
* Outil de recherche OCR\\
\\
Capacités:\\
de 5Go à 16 To, 10Go max / fichier uploadé\\
\\
Applications pour:\\
Windows, Mac, Android, iOS (prochainement)\\


\section{iCloud}

\subsection{Présentation}
Crée en 2011, cette application de la firme Apple regroupe iTunes in the Cloud, Photo stream, Calendar, Mail, et Contacts. Il permet la synchronisation de différents éléments (applications, livres, documents, sauvegardes) entre divers appareils. Il ne s'agit pas à proprement parler d'un service par navigateur, mais il possède une interface web pour contrôler certaines de ses données.\\

\subsection{Description}
DONNÉES BRUTES\\
Fonctionnalités\\
* L'interface web ne permet pas un contrôle total des fichiers (juste iWork et Photos)\\
* Le partage de fichiers n'est permis que via iWork\\
* Streaming de fichiers multimédia\\
* Intégré à iOS\\
* Backup automatique\\
\\
Capacités:\\
de 5Go à 55Go, 250Mo max / fichier uploadé\\
\\
Applications pour:\\
Windows, Mac, iOS\\


\section{Windows Live Skydrive}

\subsection{Présentation}
Il succède à Windows Live Folders distribué en 2007. Il peut s'utiliser aussi bien en tant que service web qu'en tant qu'application lourde, permettant une synchronisation avec les données locales.\\

\subsection{Description}
DONNÉES BRUTES\\
Fonctionnalités\\
* Streaming de fichiers multimédia\\
* Permet de ne synchroniser que certains dossiers préalablement sélectionnés\\
* Travail collaboratif\\
* Intégration à Office et Windows Phone\\
* Possibilité d'accès à distance aux ordinateurs possédant le client\\
\\
Capacités:\\
de 7Go à 107Go, 2Go max / fichier uploadé\\
\\
Applications pour:\\
Windows Mac, Windows Phone, Android (prochainement), iOS\\
\section{Windows Live Skydrive}
\thispagestyle{EIP} % seems mandatory

\subsection{Presentation}
Windows Live Skydrive is a solution to syncronize and share files. Developped by Microsoft it can be used with a Web Browser or with native applications on several systems.

\subsection{History}
Windows Live Folders was released to general public in August 2007. It changes its name the same month and becomes Windows Live Skydrive. In the begining it offered 5go of storage space per user, this was raised to 25Go in 2008 before beeing lowered again to 7Go in 2012.\\

\subsection{Description}
It can be used as a web service or as a rich native application, allosing for syncronisation with local data.\\
A service for colaborative working is also available for users to work on the same files.\\
It also offers distant access to computers where Windows Live Skydrive is installed.\\
\\
Its storage capacity is between 7 and 107Go, with a limit of 2go per file.\\
Included in Office and Windows Phone it has applications for Windows, Mac OS, iOS and soon for Android.\\

\subsection{Critics}
Critics towards Google Drive can also be applied to Windows Live Skydrive. Furthermore syncronisation is not possible on all plateforms which restricts the users.

\section{Owncloud}
\thispagestyle{EIP} % seems mandatory
\subsection{Présentation}
OwnCloud est une application web de stockage en ligne. Sous licence libre (AGPL), il peut-être installé sur n'importe quel serveur disposant de PHP et de SQL.\\
C'est une solution à installer soi-même, qui ne propose donc pas de louer un espace de stockage, mais d'en créer un.\\

\subsection{Historique}
Annoncé lors du Camp KDE 2010, le projet a bien évolué depuis. Le développement est très actif et suivi par la communauté (plus de 10000 commits et 1500 rapports de bogues en seulement 3 ans).\\

En 2011, une entreprise s'est créée autour du projet, proposant des services plus avancés pour les entreprises. Fin 2012, la société a réalisé une levée de fonds de 2,5 millions de dollars.

\subsection{Description}
OwnCloud repose sur un système d'applications. C'est donc un système modulaire, qui bénificie de très nombreuses fonctionnalités (édition de fichier, streaming de musique, synchronisation de calendrier et de contacts, galerie photo, etc…).\\
Plusieurs clients sont disponibles sur la plupart des plateformes, ainsi qu'une version web.\\

\subsection{Critiques}
Beaucoup d'avis d'utilisateurs critiquent les nombreux bogues d'OwnCloud, ainsi que la qualité de son code. Son plus gros avantage est qu'il permet d'installer rapidement une solution multi-usage, mais il semblerait qu'il se montre peu fiable sur le long terme.
Pour l'échange de fichiers, OwnCloud repose sur la technologie WebDAV, qui montre vite ses limites avec de gros fichiers.
À partir de la version 4.5, il est possible de s'interfacer avec Google Drive et Dropbox. Cette fonctionalité, encore expérimentale, et celle qui se rapproche le plus du Cloud Computing. En effet, OwnCloud ne permet pas de répartir la charge sur plusieurs serveurs, notion inhérente au Cloud Computing.

\section{iCloud}
\thispagestyle{EIP} % seems mandatory

\subsection{Presentation}
Mostly usefull for Apple clients, iCloud allows suncronisationm using the cloud, with different Apple products. It is provided with iOS.\\

\subsection{History}
Created in 2011, this software made by Apple combines iTunes in the Cloud, Photo stream, Calendar, Mail and Contact applications.

\subsection{Description}
This application allows for syncronisation of diferent type of files (Applications, books, documents, backups) between diferent terminals. It is not a web browser based service but it comes quite close. It provides a web interface to control some of the data.\\
Its users can easily share there files using the iWork suite and profit from automatic backup services.\\
\\
They can use up to between 5 and 55Go of space and they can store files no bigger than 250Mo.\\
It is integrated to iOS and applications for Windows and Mac OS exist.\\

\subsection{Critics}
This software is available for very few different systems and the web interface is very limited. Because it only allows acces to iWork data. To fully enjoy this product the user must only use Apple products.\\

\section{Syncany}
\thispagestyle{EIP}
\subsection{Présentation}
Syncany is a storage application that is meant to be multi-purpose. It provides the possibility to save your data on very differents supports like FTP, IMAP, SSH, CIFS, Amazon S3... The software must permit to manage ones files in a very flexible and personal way.

\subsection{Historique}
This project was launched on april 27th 2011. This project was very followed at the beggining, before loosing his attractivity.

In october 2011, this project became a student project at the Mannheim University. From this moment, its public activity strongly reduced. By now, there is almost no communication as the last official message is from april 2012.

\subsection{Description}
Syncany wants to be a thick crossplatform application which would allow secure storage (local encrypted data) and would abstract storage system. So it bears a large number of different storage ways, in a transparent maner :

\begin{itemize}
\renewcommand{\labelitemi}{$\bullet$}
\item Local Folder
\item FTP
\item IMAP
\item Google Storage
\item Amazon S3
\item Rackspace Cloud Files
\item WebDAV
\item Picasa Web Albums
\item Windows Share (NetBIOS/CIFS)
\item Box.net
\item SFTP/SSH
\end{itemize}

\vspace{1cm}

Users can with only one interface store their files in their medium of choice. It is a very flexible solution because you are not obliged to use the software itself to access, add and edit files. It allows to create various workflows, to ensure that the storage medium will always be available and even after Syncany ceases to exist, files will always be reachable.

\subsection{Critiques}
Even if Syncany´s offer seems to be very attractive, the leak of project-realted communication tends to asssume its death. And even if source code was available at the begining, it is not sufficent to be used in real situations.

As there is no official announcement of the project´s death, it is difficult to know if it is possible to contribute or even resume the work already done.

This solution is, by now, more a proof of concept than a ended software.

\chapter{Ubuntu One}
\thispagestyle{EIP} % seems mandatory
\section{Présentation}
Ubuntu One est l'equivalent libre de dropbox, avec une integration du contenu multimedia en plus. Porte par Cannonical, la societe derriere Ubuntu, cette solution fonctionne avec un client et protocole libre. Mais un serveur proprietaire.

Ubuntu One permet de synchroniser ses documents, les partager publiquement, acheter de la musique, sauvegarder ses contacts et echanger des fichiers imposants.

\section{Historique}
Ubuntu One fait sa premiere apparition publique debut 2009 et est inclu par defaut dans ubuntu depuis la version 9.10 de la distribution. Si la solution a eu du mal a convaincre les utilisateurs a ses debuts a cause du serveur proprietaire, elle s'est depuis beaucoup developpe.

Des clients multiplateforme ont notamment ete lances pour attirer plus d'utilisateurs. Le logiciel est officiellememt distribues sur :

\begin{itemize}
\renewcommand{\labelitemi}{$\bullet$}
\item Ubuntu
\item Mac Os X
\item Windows
\item Iphone
\item Android
\end{itemize}

Malgre sa provenance, la majorite des utilisateurs d'Ubuntu One sont sous windows.

Afin de developper son offre, Cannonical a aussi rajoute des fonctions multimedia a son logiciel. Il est possible d'acheter de la musique qui sera livree directement sur le compte Ubuntu One de l'utilisateur. Cette musique pourra etre diffuse directement depuis Ubuntu One.

Il est aussi possible d'utiliser Ubuntu one comme extension a Thunderbird pour parteger des documents trop lourds pour etre en piece jointe.

Malgre l'ouverture du protocole et du client, il n'existe pas encore de serveur libre pour Ubuntu one.

En Juillet 2011, Cannonical annoncait atteindre le million d'utilisateurs, toutes plateformes confondues.

\section{Description}
Le logiciel est construit autour de trois briques : Le client, le protocle et le serveur. Les deux premiers sont libres et les sources sont a disposition du public. C'est notamment grace a cela que l'extension a Thunderbird a ete possible.

\subsection{Le client}
Le client est code en python et utilise Twisted pour la gestion du reseau. L'interface graphique a evolue au cours du temps, mais elle est aujourd'hui en Qt, ce qui permet d'avoir une experience identique quelque soit la plateforme.

Une interface web est aussi disponible pour consulter ses documents.

\subsection{Le protocole}

Le protocole de communication utilise par Ubuntu One a ete cree pour l'occasion. Les equipes de developpement ont choisi de repartir de zero plutot que d'utiliser un protocole existant pour des soucis de performance.

Ce protocole est base sur les protobuffers, developpes a la base par google. Cela permet au protocole d'evoluer simplement, de limiter les problemes de serialisation, limiter la taille des donnees et de pouvoir evoluer facilement techniquement. En effet, des bibliotheques permettant d'utiliser les protobuffers sont disponibles dans de nombreux langages.

Cannonical a libere une implementation python du protocole.

\subsection{Le serveur}

Le serveur est donc l'unique brique dont nous n'avons pas tous les details, puisque nous n'avons pas les sources. Nous connaissons neanmoins certains details.

Les donnees stockees sur Ubuntu One sont envoyes sur le cloud Amazon S3. C'est notamment a cause de ce lien avec amazon que le code du serveur d'Ubuntu one ne peut etre libere par Cannonical.

Pour l'authentification, cannonical a fait le choix de OAuth. Sauf pour leur client web, ou une autre solution est utilisee.

Une API REST partielle est egalement disponible pour communiquer avec le serveur. Elle permet notamment de distribuer les documents qui ont ete partages publiquement a travers Ubuntu One.

\section{Critiques}
Le plus gros manque actuel d'Ubuntu One est l'absence de serveur libre. Le logiciel, comme la majorite de ses concurrents, ne permet pas de reprendre le controle de ses donnees.

Par rapport a ses concurrents, le service souffre aussi d'un plus faible nombre d'utilisateurs. Les fonctionnalitees de partages prives sont donc plus complique a exploite, les deux utilisateurs devant absoluement disposer d'une compte Ubuntu One.


\chapter{Conclusion}

\end{document}
